\subsection{Design reconsiderations}
\label{sec:later}

The four sections above describe each subsystem as delivered, but four of those designs are not what was built first: each replaced an earlier version once the complete chain existed and the output could be measured end to end, which is the point at which faults invisible inside a single subsystem became apparent. They are collected here, separately from the component descriptions above, because each is a change of \emph{approach} rather than a correction of code, and because the reason for each lies outside the subsystem it altered. Several of these reconsiderations are referenced from Sections~\ref{sec:writerid} and~\ref{sec:synthesis} above and are given their full justification here.

\subsubsection{A second writer classifier, judging shape instead of ink}
\label{sec:later-strokenorm}

The first writer classifier was trained on grey-scale crops and read real held-out lines at 99.78\,\% over 1358 lines, which satisfied requirement R2 on photographed pages. Measured on \emph{synthesised} output, however, the same classifier identified the intended writer in only about 28\,\% of trials. The cause is a property of the machine rather than of the network: the gantry draws every stroke with one pen at one width and one pressure, whereas real handwriting carries a thickness and an ink density that a convolutional network exploits readily when the training lines are few. A classifier that leans on those cues therefore predicts neither what the machine can draw nor how a reader recognises a hand, which is by geometry.

The response was to train a second classifier of the same architecture on stroke-normalised crops, in which the ink is thinned to a centre line and redrawn at a constant width proportional to the letter height, so thickness is removed by construction (Section~\ref{sec:wid-strokenorm}). This is the reconsideration that matters most for requirement R1: style can be judged, and reproduced, on geometry alone once pen pressure and width are taken out of the comparison. The two classifiers are used for different purposes, and the division is kept throughout the system: the ink classifier serves real photographed pages, and the shape classifier serves the candidate search of Section~\ref{sec:syn-search}. Neither replaced the other, so the writer-identification block carries two sets of weights rather than one, at the cost of double the training and storage for a capability the first classifier could not provide at all.

\subsubsection{End-to-end training in place of a frozen backbone}
\label{sec:later-endtoend}

The shape classifier was first trained by fitting only the 13\,645 parameters of its head on feature vectors extracted once from a frozen backbone. That recipe is a training fault, set out in Section~\ref{sec:writerid}: because the features were extracted once, no augmentation was re-sampled between epochs and every line contributed the same four fixed vectors for all 500 epochs. Training all 966\,445 parameters end to end instead raised the same measurement from 83.36\,\% to 93.67\,\% on the same 1358 real lines, and removed a failure in which 118 of one writer's 148 lines had been assigned to a single neighbour. Every figure quoted for the shape classifier elsewhere in this report is from the end-to-end model; the frozen-head figures describe the superseded one, and are kept here only as the evidence for why the change was made.

\subsubsection{Selecting a candidate on both criteria instead of legibility alone}
\label{sec:later-joint}

The candidate search first kept whichever of the $N$ draws the text recogniser read most accurately. That raised legibility and, for some writers, lowered writer identification, because the most legible of several random draws has no reason also to be the most characteristic. The score was therefore changed to the harmonic mean of the recogniser's character accuracy and the shape classifier's confidence in the intended writer, Equation~\ref{eq:syn-score}, dominated by whichever of the two is smaller, so a candidate cannot win on one criterion while failing the other. The repair loop was changed for the same reason, accepting a repaired line only when the joint score improves, so a repair that fixes a letter but makes the line resemble another writer is rejected rather than taken. The evidence for both changes is Table~\ref{tab:syn-evid}.

\subsubsection{Line-level style calibration}
\label{sec:later-cal}

Glyphs are stored and placed one instance at a time, so a line assembled from them does not automatically reproduce the writer's line-level slant and ascender reach. The calibration of Section~\ref{sec:syn-glyph} was added to correct that drift, measuring a rendered probe with the same functions that measured the writer's real pages and returning a shear and reach correction. The drift it corrects was itself measured before the change: writer 150's ascenders came out about 7\,\% too tall and writer 551's slant about $9^\circ$ too steep against their own real lines. Unlike the other three changes, this one has no clean before-and-after pair in Table~\ref{tab:syn-evid}, so its benefit rests on the probe agreeing with the stored reference rather than on a judge's score, and Section~\ref{sec:syn-limits} records that this agreement is partly self-fulfilling, since it is measured by the same estimator on both sides.
\mustdo{measure writer identification and character accuracy with the calibration of Section~\ref{sec:syn-glyph} disabled and enabled, over the same sentences and seeds as the other rows of Table~\ref{tab:syn-evid}, so that this change carries the same kind of evidence as the other three.}

\subsubsection{A reconsideration not yet carried out: detecting the need for an anchor automatically}

One further change was identified during the glyph-library investigation of this project but has not been implemented in the delivered system, and is recorded here as a design reconsideration still open rather than as a completed change, since the distinction between what was changed and what is merely proposed matters to how this section should be read. The anchor-substitution mechanism of Section~\ref{sec:syn-glyph} is driven entirely by the per-glyph prior distance of Equation~\ref{eq:sty-prior} and the single tuned parameter $\lambda$, applied uniformly within a writer. An examination of the glyph cutting and the cross-writer prior distribution for the most strongly joined personal writer in this project's data showed that her low connectedness-adjusted legibility is not a scattered collection of individually poor letters but a structural property of her hand: her forced-alignment cuts fall through genuine strokes far more often than for a lightly joined writer, because her connectedness statistic of 0.962 (Table~\ref{tab:sty-writers}) leaves almost no low, thin ligature for Equation~\ref{eq:sty-cut} to find. The reconsideration this suggests is to use the writer's \emph{own} connectedness and the \emph{distribution} of her glyphs' prior distances, rather than $\lambda$ alone, to decide automatically whether a writer needs anchor substitution forced on by default, instead of relying on a value tuned by hand per writer. A substitution experiment on this writer's own held-out sentences, forcing full anchor use rather than the tuned $\lambda$, raised the shape classifier's confidence in her identity from 63--77\,\% to 94\,\% and gave a modest further gain in character accuracy, which is evidence the direction is worth pursuing, though word accuracy stayed poor regardless of the glyph source, pointing to a second, separate limitation in the synthesiser's joining geometry that this reconsideration alone would not fix.
\mustdo{implement and measure the connectedness-triggered default described here across all thirteen writers, not only the one case it was explored on.}

\mustdo{confirm that these five are the complete set of reconsiderations made or identified during the project, and add any others that belong here --- for example a change to the segmentation thresholds, to the number of writers, or to the requirement allocation --- each with what was measured before and after.}
